% GENERATED FILE. Edit canonical lessons, then run npm run generate:books.
% canonical-source-hash: fnv1a64:2cbb52c16ee6d82e
% canonical-lessons: MR-C375-buddhi, MR-C375-smaransakti, MR-C375-ekagrata, MR-C375-akalan-hone, MR-C375-sarav-karne

\chapter{The Mind at Work}
\label{ch:mr-the-mind-at-work}
\hlchaptermodality{375}

\begin{chapteropening}
\textbf{By the end of this chapter:} \emph{I can talk about intellect, memory, concentration and practice.}

Complete the last lesson of chapter 375: I can talk about intellect, memory, concentration and practice.
\end{chapteropening}

\section[\texorpdfstring{\mr{बुद्धी} (buddhī)}{बुद्धी (buddhī)}]{\mr{बुद्धी} (buddhī) — intellect; intelligence}

\label{lesson:MR-C375-buddhi}



\begin{quote}
Before the new one: say the Marathi for a definition, then the Marathi for a science.
\end{quote}

\subsection*{You'll want to know: \mr{बुद्धी}}
\textbf{\mr{बुद्धी}} — \emph{buddhī} — \textquotedblleft{}intellect; intelligence\textquotedblright{}.

\begin{quote}
\mr{हा} \mr{खेळ} \mr{बुद्धीचा} \mr{आहे}, \mr{ताकदीचा} \mr{नाही}.
\end{quote}

\begin{quote}
\emph{hā kheḷ buddhīcā āhe, tākdīcā nāhī.}
\end{quote}

\begin{quote}
This game is one of brains, not strength.
\end{quote}

Feminine.

\mr{बुद्धीचा} \mr{खेळ}, a game of wits.

\begin{cousinweb}
\textbf{\mr{बुद्धी}} (\emph{buddhī}) is a tatsama, taken from Sanskrit \emph{buddhi}, \textquotedblleft{}intelligence, discernment\textquotedblright{} (CDIAL 9277).
\end{cousinweb}

\begin{grammarlens}[title={What you've built}]
One more word for giving reasons and opinions.
\end{grammarlens}

\subsection*{Your turn}
\begin{itemize}
\raggedright
  \item \emph{Say it:} \emph{buddhī}
  \item \emph{Say it:} \emph{buddhī}, once more
  \item \emph{Say it:} \emph{hā kheḷ buddhīcā āhe, tākdīcā nāhī.} — This game is one of brains, not strength.
  \item \emph{Recall:} say the Marathi for a definition, then the Marathi for a science, then say \emph{buddhī} again
\end{itemize}

\begin{tcolorbox}[breakable,colback=teal!4,colframe=teal!35!black,title={Before you move on}]
What does \mr{बुद्धी} mean? (\textquotedblleft{}Intellect; intelligence\textquotedblright{}.) Say it once more.
\end{tcolorbox}
\section[\texorpdfstring{\mr{स्मरणशक्ती} (smaraṇśaktī)}{स्मरणशक्ती (smaraṇśaktī)}]{\mr{स्मरणशक्ती} (smaraṇśaktī) — memory power; memory}

\label{lesson:MR-C375-smaransakti}



\begin{quote}
Before the new one: say the Marathi for a science, then the Marathi for intellect.
\end{quote}

\subsection*{You'll want to know: \mr{स्मरणशक्ती}}
\textbf{\mr{स्मरणशक्ती}} — \emph{smaraṇśaktī} — \textquotedblleft{}memory power; memory\textquotedblright{}.

\begin{quote}
\mr{आजीची} \mr{स्मरणशक्ती} \mr{अजूनही} \mr{चांगली} \mr{आहे}.
\end{quote}

\begin{quote}
\emph{ājīcī smaraṇśaktī ajūnhī cāṅglī āhe.}
\end{quote}

\begin{quote}
Grandmother's memory is still good.
\end{quote}

Feminine.

\mr{स्मरण} + \mr{शक्ती}, remembering + power.

\begin{cousinweb}
\textbf{\mr{स्मरणशक्ती}} (\emph{smaraṇśaktī}) is a learned compound built from Sanskrit \emph{smaraṇa}, \textquotedblleft{}remembering\textquotedblright{}, and \emph{śakti}, \textquotedblleft{}power\textquotedblright{} (Monier-Williams).
\end{cousinweb}

\begin{grammarlens}[title={What you've built}]
One more word for giving reasons and opinions.
\end{grammarlens}

\subsection*{Your turn}
\begin{itemize}
\raggedright
  \item \emph{Say it:} \emph{smaraṇśaktī}
  \item \emph{Say it:} \emph{smaraṇśaktī}, once more
  \item \emph{Say it:} \emph{ājīcī smaraṇśaktī ajūnhī cāṅglī āhe.} — Grandmother's memory is still good.
  \item \emph{Recall:} say the Marathi for a science, then the Marathi for intellect, then say \emph{smaraṇśaktī} again
\end{itemize}

\begin{tcolorbox}[breakable,colback=teal!4,colframe=teal!35!black,title={Before you move on}]
What does \mr{स्मरणशक्ती} mean? (\textquotedblleft{}Memory power; memory\textquotedblright{}.) Say it once more.
\end{tcolorbox}
\section[\texorpdfstring{\mr{एकाग्रता} (ekāgratā)}{एकाग्रता (ekāgratā)}]{\mr{एकाग्रता} (ekāgratā) — concentration}

\label{lesson:MR-C375-ekagrata}



\begin{quote}
Before the new one: say the Marathi for intellect, then the Marathi for memory power.
\end{quote}

\subsection*{You'll want to know: \mr{एकाग्रता}}
\textbf{\mr{एकाग्रता}} — \emph{ekāgratā} — \textquotedblleft{}concentration\textquotedblright{}.

\begin{quote}
\mr{अभ्यास} \mr{करताना} \mr{एकाग्रता} \mr{हवी}.
\end{quote}

\begin{quote}
\emph{abhyās kartānā ekāgratā havī.}
\end{quote}

\begin{quote}
You need concentration while studying.
\end{quote}

Feminine.

\mr{एकाग्र}, intent on one point, + -\mr{ता}.

\begin{cousinweb}
\textbf{\mr{एकाग्रता}} (\emph{ekāgratā}) is a tatsama, taken from Sanskrit \emph{ekāgratā}, \textquotedblleft{}intentness on one object\textquotedblright{} (Monier-Williams).
\end{cousinweb}

\begin{grammarlens}[title={What you've built}]
One more word for giving reasons and opinions.
\end{grammarlens}

\subsection*{Your turn}
\begin{itemize}
\raggedright
  \item \emph{Say it:} \emph{ekāgratā}
  \item \emph{Say it:} \emph{ekāgratā}, once more
  \item \emph{Say it:} \emph{abhyās kartānā ekāgratā havī.} — You need concentration while studying.
  \item \emph{Recall:} say the Marathi for intellect, then the Marathi for memory power, then say \emph{ekāgratā} again
\end{itemize}

\begin{tcolorbox}[breakable,colback=teal!4,colframe=teal!35!black,title={Before you move on}]
What does \mr{एकाग्रता} mean? (\textquotedblleft{}Concentration\textquotedblright{}.) Say it once more.
\end{tcolorbox}
\section[\texorpdfstring{\mr{आकलन} \mr{होणे} (ākalan hoṇe)}{आकलन होणे (ākalan hoṇe)}]{\mr{आकलन} \mr{होणे} (ākalan hoṇe) — to comprehend; to grasp}

\label{lesson:MR-C375-akalan-hone}



\begin{quote}
Before the new one: say the Marathi for memory power, then the Marathi for concentration.
\end{quote}

\subsection*{You'll want to know: \mr{आकलन} \mr{होणे}}
\textbf{\mr{आकलन} \mr{होणे}} — \emph{ākalan hoṇe} — \textquotedblleft{}to comprehend; to grasp\textquotedblright{}.

\begin{quote}
\mr{पुन्हा} \mr{वाचल्यावर} \mr{मला} \mr{धड्याचं} \mr{नीट} \mr{आकलन} \mr{झालं}.
\end{quote}

\begin{quote}
\emph{punhā vācalyāvar malā dhaḍyācaṁ nīṭ ākalan jhālaṁ.}
\end{quote}

\begin{quote}
After reading it again, I grasped the lesson properly.
\end{quote}

Past (to him …): \mr{त्याला} \mr{आकलन} \mr{झालं}, \emph{tyālā ākalan jhālaṁ}.

Dative subject: \mr{मला} \mr{आकलन} \mr{होतं}; what is grasped takes -\mr{चं}.

\begin{cousinweb}
\textbf{\mr{आकलन} \mr{होणे}} (\emph{ākalan hoṇe}) joins the tatsama \mr{आकलन}, Sanskrit \emph{ākalana}, \textquotedblleft{}fastening, reckoning\textquotedblright{} (Monier-Williams), to \mr{होणे}, \textquotedblleft{}to become\textquotedblright{}; Marathi uses \mr{आकलन} for taking in what something means.
\end{cousinweb}

\begin{grammarlens}[title={What you've built}]
One more word for giving reasons and opinions.
\end{grammarlens}

\subsection*{Your turn}
\begin{itemize}
\raggedright
  \item \emph{Say it:} \emph{ākalan hoṇe}
  \item \emph{Say it:} \emph{ākalan hoṇe}, once more
  \item \emph{Say it:} \emph{punhā vācalyāvar malā dhaḍyācaṁ nīṭ ākalan jhālaṁ.} — After reading it again, I grasped the lesson properly.
  \item \emph{Recall:} say the Marathi for memory power, then the Marathi for concentration, then say \emph{ākalan hoṇe} again
\end{itemize}

\begin{tcolorbox}[breakable,colback=teal!4,colframe=teal!35!black,title={Before you move on}]
What does \mr{आकलन} \mr{होणे} mean? (\textquotedblleft{}To comprehend; to grasp\textquotedblright{}.) Say it once more.
\end{tcolorbox}
\section[\texorpdfstring{\mr{सराव} \mr{करणे} (sarāv karṇe)}{सराव करणे (sarāv karṇe)}]{\mr{सराव} \mr{करणे} (sarāv karṇe) — to practise}

\label{lesson:MR-C375-sarav-karne}



\begin{quote}
Before the new one: say the Marathi for concentration, then the Marathi for to comprehend.
\end{quote}

\subsection*{You'll want to know: \mr{सराव} \mr{करणे}}
\textbf{\mr{सराव} \mr{करणे}} — \emph{sarāv karṇe} — \textquotedblleft{}to practise\textquotedblright{}.

\begin{quote}
\mr{ती} \mr{रोज} \mr{एक} \mr{तास} \mr{गाण्याचा} \mr{सराव} \mr{करते}.
\end{quote}

\begin{quote}
\emph{tī roj ek tās gāṇyācā sarāv karte.}
\end{quote}

\begin{quote}
She practises singing for an hour every day.
\end{quote}

Past, with -\mr{ने} (he …): \mr{त्याने} \mr{सराव} \mr{केला}, \emph{tyāne sarāv kelā}.

What is practised takes -\mr{चा}: \mr{गाण्याचा} \mr{सराव}.

\begin{cousinweb}
\textbf{\mr{सराव} \mr{करणे}} (\emph{sarāv karṇe}) joins \mr{सराव}, \textquotedblleft{}practice\textquotedblright{}, a Marathi noun whose origin is not settled, to \mr{करणे}, \textquotedblleft{}to do\textquotedblright{}.
\end{cousinweb}

\begin{grammarlens}[title={What you've built}]
One more word for giving reasons and opinions.
\end{grammarlens}

\subsection*{Your turn}
\begin{itemize}
\raggedright
  \item \emph{Say it:} \emph{sarāv karṇe}
  \item \emph{Say it:} \emph{sarāv karṇe}, once more
  \item \emph{Say it:} \emph{tī roj ek tās gāṇyācā sarāv karte.} — She practises singing for an hour every day.
  \item \emph{Recall:} say the Marathi for concentration, then the Marathi for to comprehend, then say \emph{sarāv karṇe} again
\end{itemize}

\begin{tcolorbox}[breakable,colback=teal!4,colframe=teal!35!black,title={Before you move on}]
What does \mr{सराव} \mr{करणे} mean? (\textquotedblleft{}To practise\textquotedblright{}.) Say it once more.
\end{tcolorbox}
